\documentclass[11pt,a4paper]{article}

%% ── Packages ──────────────────────────────────────────────────────────────
\usepackage[utf8]{inputenc}
\usepackage[T1]{fontenc}
\usepackage[spanish]{babel}
\usepackage[top=2.5cm,bottom=2.5cm,left=2.5cm,right=2.5cm]{geometry}
\usepackage{amsmath,amssymb}
\usepackage{xcolor}
\usepackage[most,breakable]{tcolorbox}
\usepackage{enumitem}
\usepackage{booktabs}
\usepackage{array}
\usepackage{fancyhdr}
\usepackage[hidelinks,colorlinks=true,linkcolor=itbaBlue,urlcolor=itbaBlue]{hyperref}
\usepackage{lmodern}

%% ── Colors ────────────────────────────────────────────────────────────────
\definecolor{itbaAmber}{HTML}{F2A32B}
\definecolor{itbaBlue}{HTML}{3FA8F5}
\definecolor{itbaRed}{HTML}{F25C54}
\definecolor{itbaGreen}{HTML}{3DD68C}
\definecolor{darkText}{HTML}{1A2535}
\definecolor{lightGray}{HTML}{F5F6FA}
\definecolor{medGray}{HTML}{DCDFE8}

%% ── tcolorbox styles ──────────────────────────────────────────────────────
% Definition / Theorem — blue
\tcbset{
  defstyle/.style={
    enhanced, breakable,
    colback=itbaBlue!8, colframe=itbaBlue, coltitle=white,
    fonttitle=\bfseries\small, arc=4pt,
    boxrule=1pt, left=6pt, right=6pt, top=4pt, bottom=4pt,
    attach boxed title to top left={yshift=-2mm,xshift=6mm},
    boxed title style={colback=itbaBlue, arc=3pt, boxrule=0pt}
  },
  % Step-by-step — amber
  methodstyle/.style={
    enhanced, breakable,
    colback=itbaAmber!10, colframe=itbaAmber, coltitle=darkText,
    fonttitle=\bfseries\small, arc=4pt,
    boxrule=1.5pt, left=6pt, right=6pt, top=4pt, bottom=4pt,
    attach boxed title to top left={yshift=-2mm,xshift=6mm},
    boxed title style={colback=itbaAmber, arc=3pt, boxrule=0pt}
  },
  % Example header
  examplestyle/.style={
    enhanced,
    colback=darkText!5, colframe=darkText!40, coltitle=white,
    fonttitle=\bfseries\small, arc=4pt,
    boxrule=1pt, left=6pt, right=6pt, top=4pt, bottom=4pt,
    attach boxed title to top left={yshift=-2mm,xshift=6mm},
    boxed title style={colback=darkText!80, arc=3pt, boxrule=0pt}
  },
  % Conclusion SÍ — green
  yestyle/.style={
    enhanced,
    colback=itbaGreen!12, colframe=itbaGreen, coltitle=darkText,
    fonttitle=\bfseries\small, arc=4pt,
    boxrule=1.2pt, left=6pt, right=6pt, top=4pt, bottom=4pt,
    attach boxed title to top left={yshift=-2mm,xshift=6mm},
    boxed title style={colback=itbaGreen, arc=3pt, boxrule=0pt}
  },
  % Conclusion NO — red
  nostyle/.style={
    enhanced,
    colback=itbaRed!10, colframe=itbaRed, coltitle=white,
    fonttitle=\bfseries\small, arc=4pt,
    boxrule=1.2pt, left=6pt, right=6pt, top=4pt, bottom=4pt,
    attach boxed title to top left={yshift=-2mm,xshift=6mm},
    boxed title style={colback=itbaRed, arc=3pt, boxrule=0pt}
  },
  % Memory hook — amber left border
  memostyle/.style={
    enhanced,
    colback=itbaAmber!8, colframe=itbaAmber,
    leftrule=4pt, toprule=0pt, bottomrule=0pt, rightrule=0pt,
    arc=0pt, left=8pt, right=6pt, top=4pt, bottom=4pt,
    fontupper=\small\itshape
  }
}

%% ── Header / Footer ───────────────────────────────────────────────────────
\pagestyle{fancy}
\fancyhf{}
\renewcommand{\headrulewidth}{0.6pt}
\renewcommand{\footrulewidth}{0.4pt}
\fancyhead[L]{\small\color{darkText!70}\textbf{ITBA} · Criptografía}
\fancyhead[R]{\small\color{darkText!70}Secreto Perfecto}
\fancyfoot[C]{\small\color{darkText!60}\thepage}
\fancyfoot[R]{\small\color{darkText!50}Primer Parte}

%% ── Document ──────────────────────────────────────────────────────────────
\begin{document}

\color{darkText}

%% Title
\begin{tcolorbox}[
  enhanced, arc=6pt, boxrule=0pt,
  colback=darkText, coltext=white,
  left=12pt, right=12pt, top=10pt, bottom=10pt
]
  {\Large\bfseries Secreto Perfecto (Perfect Secrecy)}\\[4pt]
  {\small Criptografía · ITBA · Resumen de estudio}
\end{tcolorbox}

\vspace{6pt}
\tableofcontents
\vspace{10pt}

%% ══════════════════════════════════════════════════════════════════════════
\section{Definición y definición alternativa}
%% ══════════════════════════════════════════════════════════════════════════

\begin{tcolorbox}[defstyle, title={Definición: Secreto perfecto (Shannon)}]
Un criptosistema $(\mathsf{Gen}, \mathsf{Enc}, \mathsf{Dec})$ tiene \textbf{secreto perfecto} si para toda distribución sobre $\mathcal{M}$, para todo mensaje $m$ y todo cifrado $c$ con $\Pr[C{=}c]>0$:
\[
  \Pr[M{=}m \mid C{=}c] \;=\; \Pr[M{=}m]
\]
\end{tcolorbox}

\begin{tcolorbox}[memostyle]
\textbf{Intuición:} conocer el cifrado $c$ \emph{no} aporta ninguna información sobre el mensaje.
\end{tcolorbox}

\medskip
\begin{tcolorbox}[defstyle, title={Definición alternativa (equivalente, por encripción)}]
Para todo par de mensajes $m, m'$ y todo cifrado $c$:
\[
  \Pr[\mathsf{Enc}_k(m) = c] \;=\; \Pr[\mathsf{Enc}_k(m') = c]
\]
El cifrado de \emph{cualquier} mensaje tiene exactamente la misma distribución de probabilidad.
\end{tcolorbox}

\medskip
\begin{tcolorbox}[memostyle]
\textbf{La paradoja:} $C$ y $M$ están relacionados por una función ($c = \mathsf{Enc}_k(m)$), pero la condición dice que son estadísticamente independientes. La clave aleatoria introduce independencia a pesar de la relación funcional.
\end{tcolorbox}

%% ══════════════════════════════════════════════════════════════════════════
\section{Teorema condición necesaria: $|\mathcal{K}| \geq |\mathcal{M}|$}
%% ══════════════════════════════════════════════════════════════════════════

\begin{tcolorbox}[defstyle, title={Teorema (condición necesaria)}]
\[
  \text{Secreto perfecto} \;\Longrightarrow\; |\mathcal{K}| \;\geq\; |\mathcal{C}| \;\geq\; |\mathcal{M}|
\]
Si hay menos claves que mensajes posibles, \textbf{NO puede haber secreto perfecto}.
\end{tcolorbox}

\medskip
\noindent\textbf{Resultados adicionales:}
\begin{itemize}[leftmargin=1.8em]
  \item Todo criptosistema con secreto perfecto es reducible al OTP (isomorfismo).
  \item \textbf{Consecuencia práctica:} secreto perfecto es impráctico ($|\mathcal{K}| \geq |\mathcal{M}|$ siempre).
\end{itemize}

%% ══════════════════════════════════════════════════════════════════════════
\section{Teorema de Shannon (cuando $|\mathcal{K}| = |\mathcal{M}|$)}
%% ══════════════════════════════════════════════════════════════════════════

\begin{tcolorbox}[defstyle, title={Teorema de Shannon}]
Cuando además $|\mathcal{K}| = |\mathcal{M}|$, el secreto perfecto exige:
\begin{enumerate}[label=\textbf{\arabic*.}, leftmargin=2em]
  \item Las claves se eligen con distribución \textbf{uniforme}:
    $\Pr[K{=}k] = \tfrac{1}{|\mathcal{K}|}$ para todo $k$.
  \item Para cada par $(m, c)$, existe una \textbf{única} clave $k$ tal que $\mathsf{Enc}_k(m) = c$.
\end{enumerate}
\end{tcolorbox}

%% ══════════════════════════════════════════════════════════════════════════
\section{Lema auxiliar: cómo calcular $\Pr[C{=}c]$}
%% ══════════════════════════════════════════════════════════════════════════

\begin{tcolorbox}[defstyle, title={Lema auxiliar}]
\[
  \Pr[C{=}c] \;=\; \sum_k \Pr[C{=}c \cap K{=}k]
             \;=\; \sum_k \Pr\!\bigl[M = \mathsf{Dec}_k(c)\bigr] \cdot \Pr[K{=}k]
\]
La segunda igualdad usa la independencia entre $M$ y $K$.
\end{tcolorbox}

%% ══════════════════════════════════════════════════════════════════════════
\section{Método paso a paso para demostrar o refutar secreto perfecto}
%% ══════════════════════════════════════════════════════════════════════════

\begin{tcolorbox}[methodstyle, title={Método paso a paso}]
\begin{enumerate}[label=\textbf{Paso \arabic{*} ---}, leftmargin=3.5em, itemsep=4pt]
  \item \textbf{Calcular cardinales $|\mathcal{M}|$, $|\mathcal{K}|$, $|\mathcal{C}|$}\\
    Si $|\mathcal{K}| < |\mathcal{M}|$: concluir directamente que \textbf{NO} hay secreto perfecto.

  \item \textbf{Enunciar qué definición se va a usar}
    \begin{itemize}[leftmargin=1.5em, itemsep=1pt]
      \item \textbf{Opción A (Shannon):} verificar $\Pr[M{=}m \mid C{=}c] = \Pr[M{=}m]$
      \item \textbf{Opción B (por encripción):} verificar $\Pr[\mathsf{Enc}_k(m){=}c] = \Pr[\mathsf{Enc}_k(m'){=}c]$
    \end{itemize}

  \item \textbf{Verificar el Teorema de Shannon (si $|\mathcal{K}|{=}|\mathcal{M}|$)}
    \begin{itemize}[leftmargin=1.5em, itemsep=1pt]
      \item ¿Las claves son uniformes?
      \item ¿Para cada par $(m,c)$ existe una única $k$ tal que $\mathsf{Enc}_k(m){=}c$?
    \end{itemize}

  \item \textbf{Calcular (si es necesario) usando el lema}
    \begin{itemize}[leftmargin=1.5em, itemsep=1pt]
      \item Calcular $\Pr[C{=}c]$ sumando sobre claves.
      \item Calcular $\Pr[M{=}m \mid C{=}c]$ usando Bayes.
      \item Comparar con $\Pr[M{=}m]$.
    \end{itemize}

  \item \textbf{Concluir}
    \begin{itemize}[leftmargin=1.5em, itemsep=1pt]
      \item Si todo se cumple: \textbf{SÍ} hay secreto perfecto.
      \item Si falla algo: exhibir el contraejemplo.
    \end{itemize}
\end{enumerate}
\end{tcolorbox}

%% ══════════════════════════════════════════════════════════════════════════
\section{Ejemplo 1: OTP --- tiene secreto perfecto}
%% ══════════════════════════════════════════════════════════════════════════

\begin{tcolorbox}[examplestyle, title={Ejemplo 1: One-Time Pad (OTP)}]
$\mathsf{Enc}_k(m) = m \oplus k$, clave uniforme, $|k| = |m| = n$ bits,
$|\mathcal{K}| = |\mathcal{M}| = |\mathcal{C}| = 2^n = N$.

\medskip
\textbf{Objetivo:} probar $\Pr[M{=}m \mid C{=}c] = \Pr[M{=}m]$.

\medskip
\textbf{Lema 1 --- $\Pr[C{=}c] = 1/N$:}
\begin{align*}
  \Pr[C{=}c]
    &= \sum_k \Pr[C{=}c \cap K{=}k]\\
    &= \sum_k \Pr[M = c \oplus k] \cdot \Pr[K{=}k]\\
    &= \sum_k \Pr[M = c \oplus k] \cdot \tfrac{1}{N}\\
    &= \tfrac{1}{N} \cdot \sum_k \Pr[M = c \oplus k]
     = \tfrac{1}{N} \cdot 1
     = \frac{1}{N}
\end{align*}
(Al recorrer todas las claves $k$, el valor $c \oplus k$ recorre todos los mensajes sin repetir.)

\medskip
\textbf{Parte 2 --- Conclusión:}
\begin{align*}
  \Pr[M{=}m \mid C{=}c] \cdot \Pr[C{=}c]
    &= \Pr[C{=}c \cap M{=}m]\\
    &= \Pr[K{=}c \oplus m \cap M{=}m]\\
    &= \Pr[K{=}c \oplus m] \cdot \Pr[M{=}m]
     = \tfrac{1}{N} \cdot \Pr[M{=}m]
\end{align*}
Como $\Pr[C{=}c] = 1/N$, se cancelan:
$\Pr[M{=}m \mid C{=}c] = \Pr[M{=}m]$.
\end{tcolorbox}

\begin{tcolorbox}[yestyle, title={Conclusi\'on --- S\'I hay secreto perfecto}]
El OTP con clave uniforme tiene secreto perfecto. $\blacksquare$
\end{tcolorbox}

%% ══════════════════════════════════════════════════════════════════════════
\section{Ejemplo 2: OTP con clave sesgada --- no tiene secreto perfecto}
%% ══════════════════════════════════════════════════════════════════════════

\begin{tcolorbox}[examplestyle, title={Ejemplo 2: OTP con clave no uniforme (2 bits)}]
Clave con distribución sesgada:
$\Pr[k{=}00]{=}0.3$, $\Pr[k{=}01]{=}0.1$, $\Pr[k{=}10]{=}0.4$, $\Pr[k{=}11]{=}0.2$.

\smallskip
Mensajes: $\Pr[m{=}00]{=}0.60$, $\Pr[m{=}01]{=}0.15$, $\Pr[m{=}10]{=}0.10$, $\Pr[m{=}11]{=}0.15$.

\medskip
Se intercepta $c = 01$. Calcular $\Pr[M{=}00 \mid C{=}01]$:

\medskip
\textbf{Paso (a) --- Calcular $\Pr[C{=}01]$:}

Para cada mensaje $m$, la clave que produce $01$ es $k = m \oplus 01$:
\begin{align*}
  \Pr[C{=}01]
    &= \Pr[M{=}01]\cdot\Pr[K{=}00]
     + \Pr[M{=}00]\cdot\Pr[K{=}01]
     + \Pr[M{=}11]\cdot\Pr[K{=}10]
     + \Pr[M{=}10]\cdot\Pr[K{=}11]\\
    &= 0.15 \cdot 0.3 + 0.60 \cdot 0.1 + 0.15 \cdot 0.4 + 0.10 \cdot 0.2\\
    &= 0.045 + 0.06 + 0.06 + 0.02 \;=\; 0.185
\end{align*}

\medskip
\textbf{Paso (b) --- Calcular $\Pr[C{=}01 \mid M{=}00]$:}

Si $m = 00$, la única $k$ que produce $c = 01$ es $k = 01$.
\[
  \Pr[C{=}01 \mid M{=}00] = \Pr[K{=}01] = 0.1
\]

\medskip
\textbf{Paso (c) --- Aplicar Bayes:}
\[
  \Pr[M{=}00 \mid C{=}01]
    = \frac{\Pr[C{=}01 \mid M{=}00]\cdot\Pr[M{=}00]}{\Pr[C{=}01]}
    = \frac{0.1 \times 0.60}{0.185}
    \approx 0.32
\]
\end{tcolorbox}

\begin{tcolorbox}[nostyle, title={Conclusi\'on --- NO hay secreto perfecto}]
$0.32 \neq 0.60 = \Pr[M{=}00]$. La clave sesgada rompe el secreto perfecto.
\end{tcolorbox}

%% ══════════════════════════════════════════════════════════════════════════
\section{Ejemplo 3: Cifrado de transposición con $\mathcal{K}=\{2,3\}$}
%% ══════════════════════════════════════════════════════════════════════════

\begin{tcolorbox}[examplestyle, title={Ejemplo 3: Transposición sobre alfabeto de 6 letras}]
Mensajes de 6 letras del alfabeto $\{a,b,c,d,e,f\}$:
$|\mathcal{M}| = 6^6 = 46656$. \quad Clave: $\mathcal{K} = \{2, 3\}$, $|\mathcal{K}| = 2$.
\end{tcolorbox}

\begin{tcolorbox}[nostyle, title={Conclusi\'on --- NO hay secreto perfecto}]
$|\mathcal{K}| = 2 < 46656 = |\mathcal{M}|$. Por el teorema necesario, NO hay secreto perfecto.

\smallskip
\textbf{Caso especial:} si $\mathcal{M}$ solo contiene mensajes con letras distintas,
$|\mathcal{M}| = 6! = 720$. Sigue siendo $|\mathcal{K}| = 2 < 720$. Tampoco hay secreto perfecto.
\end{tcolorbox}

%% ══════════════════════════════════════════════════════════════════════════
\section{Ejemplo 4: Cifrado af\'in --- verificar Teorema de Shannon}
%% ══════════════════════════════════════════════════════════════════════════

\begin{tcolorbox}[examplestyle, title={Ejemplo 4: Cifrado afín sobre $\mathbb{Z}_{26}$}]
Esquema: $m \in \mathbb{Z}_{26}$, $k \in \mathbb{Z}_{26}$, $c = (m + k) \bmod 26$.

\smallskip
$|\mathcal{M}| = 26$, $|\mathcal{K}| = 26$, $|\mathcal{C}| = 26$.
Se cumple $|\mathcal{K}| = |\mathcal{M}| = |\mathcal{C}|$.

\medskip
\textbf{Condici\'on 1 --- claves uniformes:}
$\mathsf{Gen}$ elige $k$ uniformemente en $\mathbb{Z}_{26}$. \checkmark

\medskip
\textbf{Condici\'on 2 --- \'unica clave para cada par $(m,c)$:}
$k = (c - m) \bmod 26$. Existe y es única. \checkmark
\end{tcolorbox}

\begin{tcolorbox}[yestyle, title={Conclusi\'on --- S\'I hay secreto perfecto}]
El cifrado af\'in sobre $\mathbb{Z}_{26}$ con clave uniforme cumple ambas condiciones
del Teorema de Shannon. $\blacksquare$
\end{tcolorbox}

%% ══════════════════════════════════════════════════════════════════════════
\section{Limitaciones del OTP}
%% ══════════════════════════════════════════════════════════════════════════

\begin{tcolorbox}[nostyle, title={Limitaciones prácticas del OTP}]
\begin{enumerate}[label=\textbf{\arabic*.}, leftmargin=2em, itemsep=4pt]
  \item \textbf{Clave tan larga como el mensaje.}
    Para un video 4K se necesita una clave del mismo tamaño.

  \item \textbf{No reutilizar la clave.}
    Si $c_1 = m_1 \oplus k$ y $c_2 = m_2 \oplus k$, entonces
    $c_1 \oplus c_2 = m_1 \oplus m_2$, filtrando información.

  \item \textbf{Clave verdaderamente aleatoria.}
    Los generadores pseudoaleatorios son deterministas y no son aceptables.
\end{enumerate}
\end{tcolorbox}

\begin{tcolorbox}[memostyle]
\textbf{Para recordar:} el OTP es información-teóricamente seguro pero logísticamente inviable para uso general. Por eso se recurre a la seguridad computacional.
\end{tcolorbox}

%% ══════════════════════════════════════════════════════════════════════════
\section{Comparación: secreto perfecto vs.\ seguridad computacional}
%% ══════════════════════════════════════════════════════════════════════════

\begin{center}
\renewcommand{\arraystretch}{1.4}
\begin{tabular}{>{\bfseries}lp{4.5cm}p{5cm}}
  \toprule
  \textbf{Criterio} & \textbf{Secreto perfecto} & \textbf{Seguridad computacional} \\
  \midrule
  Adversario      & Sin l\'imite computacional   & PPT (polytime) \\
  $|\mathcal{K}|$ vs $|\mathcal{M}|$ & $|\mathcal{K}| \geq |\mathcal{M}|$ siempre & $|\mathcal{K}| \ll |\mathcal{M}|$ posible \\
  Garant\'ia     & Incondicional (matem\'atica) & Probabil\'istica (ventaja negligible) \\
  Ejemplo         & OTP                        & AES-CTR, AES-CBC \\
  Pr\'actico      & No                         & S\'i \\
  \bottomrule
\end{tabular}
\end{center}

\medskip
\begin{tcolorbox}[memostyle]
\textbf{Para recordar:} la seguridad computacional relaja la condición del adversario. En lugar de ``ningún adversario'', se exige ``ningún adversario eficiente'' (PPT), lo que permite claves mucho más cortas que el mensaje.
\end{tcolorbox}

\end{document}
